% ============================================================
%  智能发票识别与自动报销填单系统 — 课程大作业报告
%  编译方式: XeLaTeX (Overleaf: Menu → Compiler → XeLaTeX)
% ============================================================
\documentclass[12pt,a4paper]{ctexart}

% ---------- 宏包 ----------
\usepackage{geometry}
\geometry{left=2.5cm, right=2.5cm, top=2.5cm, bottom=2.5cm}
\usepackage{graphicx}
\usepackage{float}
\usepackage{enumitem}
\usepackage{booktabs}
\usepackage{longtable}
\usepackage{xcolor}
\usepackage{fontspec}
\setmonofont{Courier New}
\usepackage{listings}
\usepackage{hyperref}

% ---------- 超链接设置 ----------
\hypersetup{
  colorlinks=true,
  linkcolor=blue,
  urlcolor=cyan,
}

% ---------- 代码块样式 ----------
\definecolor{codebg}{RGB}{245,245,245}
\lstset{
  backgroundcolor=\color{codebg},
  basicstyle=\small\ttfamily,
  breaklines=true,
  frame=single,
  framesep=4pt,
  rulecolor=\color{gray!30},
  aboveskip=6pt,
  belowskip=6pt,
  extendedchars=false,
}

% ---------- 标题深度 ----------
\setcounter{secnumdepth}{3}
\setcounter{tocdepth}{2}

% ---------- 个人信息 ----------
\newcommand{\studentname}{张文萱}
\newcommand{\studentid}{202310232056}
\newcommand{\studentmajor}{自动化}
\newcommand{\studentclass}{自动化232}
\newcommand{\submitdate}{\today}

\begin{document}

% ============================================================
%  封面
% ============================================================
\begin{titlepage}
  \centering
  \vspace*{2cm}

  {\Huge\bfseries 智能发票识别与自动报销填单系统\par}
  \vspace{0.8cm}
  {\Large —— 基于 YOLO + PaddleOCR + 大模型\par}
  \vspace{0.4cm}
  {\large （含二维码校验、重复报销检测）\par}

  \vspace{3cm}

  {\large
  \begin{tabular}{rl}
    \textbf{课程名称} & 深度学习理论及应用实践 \\
    \textbf{姓\qquad 名} & \studentname \\
    \textbf{学\qquad 号} & \studentid \\
    \textbf{专\qquad 业} & \studentmajor \\
    \textbf{班\qquad 级} & \studentclass \\
    \textbf{提交日期} & \submitdate \\
  \end{tabular}
  }

  \vfill
\end{titlepage}

% ============================================================
%  摘要
% ============================================================
\begin{abstract}
\noindent
本系统面向增值税发票报销场景，实现了一个端到端的智能发票识别与自动报销填单 Web 应用。系统以 YOLOv8 微调模型定位并裁剪发票二维码，通过 pyzbar 解析二维码中的发票号码、金额、日期等信息；使用 PaddleOCR 离线引擎对完整发票图片进行文字识别；调用多模态大模型（默认通义千问 Qwen-VL）从 OCR 结果和图片视觉布局中提取发票代码、号码、日期、购买方、销售方、金额、税额、价税合计等 7 个关键字段；将二维码解析结果与大模型提取结果进行双信息校验，不一致则阻塞提交并标红差异；校验通过后自动填入报销表单，用户补填报销人、事由、部门即可提交。提交时自动检测发票号码是否重复报销，并提供 Excel 和 PDF 双格式导出。系统采用 FastAPI + Jinja2 构建 Web 前端，SQLite 存储数据，所有处理进度通过 SSE 实时推送到前端。经过 1187 张增强数据训练，YOLO 模型 mAP50 达 0.985；端到端测试覆盖上传、检测、解析、OCR、大模型提取、校验、填单、查重、导出的完整链路，全部通过。
\end{abstract}

\newpage

% ============================================================
%  目录
% ============================================================
\tableofcontents
\newpage

% ============================================================
%  第 1 章  项目背景与意义
% ============================================================
\section{项目背景与意义}

\subsection{电子发票报销痛点}

增值税发票（含增值税专用发票、增值税普通发票及电子发票）是企业财务报销中最常见的票据类型之一。当前财务人员在处理发票报销时面临以下问题：

\begin{enumerate}[label=(\arabic*)]
  \item \textbf{手工录入效率低}：每张发票需逐字段录入发票代码、号码、日期、购买方、销售方、金额、税额、价税合计等 7 个关键字段，耗时且易出错。
  \item \textbf{重复报销风险}：同一张发票可能被不同人员或同一人员多次提交报销，缺乏自动查重机制。
  \item \textbf{数据校验困难}：发票版面信息（OCR 可读）与二维码编码信息（机器可解析）之间缺乏自动比对，人工核对不现实。
  \item \textbf{格式多样性}：纸质发票拍照、电子发票截图、PDF 电子发票等输入格式各异，需要系统具备统一的前处理能力。
\end{enumerate}

\subsection{本系统解决的问题}

本系统针对上述痛点，设计并实现了一条完整的自动化处理流水线：

\begin{itemize}
  \item \textbf{自动检测与解析}：YOLO 定位二维码区域 → pyzbar 解码 → 获取发票号码、金额、日期等结构化数据。
  \item \textbf{离线 OCR + 在线大模型}：PaddleOCR 离线识别全部文字，多模态大模型结合视觉布局和文字提取 7 个字段，兼顾隐私与智能。
  \item \textbf{双信息校验}：逐字段比对二维码数据与 OCR+大模型提取结果，全数匹配方可通过，不一致则阻塞并可视化差异。
  \item \textbf{重复检测}：按发票号码查重，已报销发票再次提交时报警并阻止入库。
  \item \textbf{自动填单与导出}：校验通过后自动填入报销表单，支持 Excel 和 PDF 双格式一键导出。
\end{itemize}


% ============================================================
%  第 2 章  相关技术与原理
% ============================================================
\section{相关技术与原理}

\subsection{YOLO 目标检测}

YOLO（You Only Look Once）是一种单阶段目标检测算法，将目标检测建模为回归问题，直接从图像像素预测边界框和类别概率。本系统使用 Ultralytics YOLOv8n（Nano 版本，参数量约 3.2M），在预训练权重基础上对增值税发票的二维码区域进行迁移微调。

YOLOv8 采用 Anchor-Free 检测头和解耦头结构（分类与回归分叉），训练时使用 CIoU 损失函数。由于发票二维码通常为正方形且位置固定（左上角或右上角），检测任务相对简单，Nano 模型即可满足精度和速度要求——单张推理仅需约 2.8ms（RTX 4060）。

\subsection{PaddleOCR 文字识别}

PaddleOCR 是百度开源的端到端 OCR 工具包，内置文本检测（DB 算法）和文本识别（CRNN + CTC）两个阶段的推理模型。本系统使用 PaddleOCR 的离线推理能力，对整张发票图片进行中英文混合文字识别，输出文本行列表及其置信度。

选择 PaddleOCR 的原因：(1) 中文识别精度在开源方案中领先；(2) 完全离线运行，不依赖网络，不向外部传输发票数据；(3) 支持 GPU 加速推理。

实际使用中对输入图片做了最长边 1024px 的缩放优化，将 OCR 耗时从 14.6s 降至约 9.4s。

\subsection{多模态大模型信息抽取}

传统方案中，从 OCR 文本提取结构化字段通常依赖正则表达式或模板匹配，对版面变化敏感且泛化性差。本系统采用多模态大模型方案：同时将发票图片（Base64 编码）和 OCR 识别文字发送给大模型，模型结合视觉布局信息和文字内容，按 JSON Schema 输出结构化提取结果。

多模态输入的关键优势：(1) 视觉布局信息帮助模型区分"购买方"和"销售方"（二者位置不同但文字相似）；(2) 表格结构的视觉理解能更准确提取金额、税额等数值字段；(3) 中文语义理解能处理换行、断句等问题。

系统设计了多后端适配器，支持通义千问（DashScope）、GPT-4o（OpenAI）、Claude（Anthropic）和 DeepSeek 四种 API，通过配置文件切换。

当大模型 API 不可用时，系统自动降级为正则规则提取器，基于增值税发票的固定关键词模式（如"发票号码："后跟 20 位数字）匹配字段。

\subsection{二维码编解码}

增值税发票的二维码遵循国家税务总局统一格式。纸质发票和电子发票的编码格式不同：

\begin{itemize}
  \item \textbf{纸质发票格式}：以逗号分隔，包含发票代码、发票号码、不含税金额、开票日期、校验码等字段。
  \item \textbf{电子发票格式}：在纸质发票基础上有额外的前导字段（票据类型标识），发票号码位于第 4 个字段、价税合计位于第 5 个字段。
\end{itemize}

本系统使用 pyzbar 库解析二维码图像，自动检测两种格式并正确提取字段。

\subsection{数据库查重逻辑}

重复报销检测是财务系统的刚需。本系统以发票号码为唯一键进行查重：同一发票号码在数据库中已存在时，提交被阻塞，前端跳转到重复报警页展示原有报销记录。

选择发票号码而非发票代码+号码组合作为查重键的原因是：表单字段可编辑，若使用组合键，用户可通过修改发票代码绕过查重。发票号码在增值税发票体系中具有唯一性，单独使用不会产生误报。


% ============================================================
%  第 3 章  系统总体设计
% ============================================================
\section{系统总体设计}

\subsection{系统架构}

系统采用经典的 Web 三层架构，但将核心业务逻辑抽离为独立的 pipeline 模块，实现业务与 Web 层的彻底解耦：

\begin{figure}[H]
\centering
\includegraphics[width=0.2\textwidth]{screenshots/architecture.drawio.png}
\caption{系统架构图}
\label{fig:architecture}
\end{figure}

\begin{itemize}
  \item \textbf{表示层（Web）}：FastAPI 路由 + 6 个 Jinja2 模板页面，负责参数绑定、结果转发和页面渲染。所有 UI 逻辑仅做展示，不含业务代码。
  \item \textbf{业务层（Pipeline）}：7 个独立模块——YOLO 检测器、QR 解码器、PaddleOCR 引擎、LLM 提取器、规则提取器、校验器、重复检测器。每个模块输入/输出清晰，可独立单元测试。
  \item \textbf{数据层（DB + Export）}：SQLAlchemy ORM + SQLite，存储发票和报销记录；Excel/PDF 导出模块接收字典数据返回文件流。
\end{itemize}

\subsection{模块划分}

\begin{longtable}{p{3cm} p{6cm} p{6cm}}
\toprule
\textbf{模块} & \textbf{文件} & \textbf{职责} \\
\midrule
YOLO 检测器 & pipeline/yolo\_detector.py & 加载预训练权重，检测并裁剪二维码区域 \\
QR 解码器   & pipeline/qr\_decoder.py   & 解析二维码，返回发票号码/金额/日期 \\
PaddleOCR   & pipeline/paddle\_ocr.py   & 离线 OCR 推理，返回文本行列表 \\
LLM 提取器  & pipeline/llm\_extractor.py & 多模态大模型提取 7 个发票字段（多后端） \\
规则提取器  & pipeline/rule\_extractor.py & 正则降级方案，大模型不可用时自动切换 \\
校验器     & pipeline/verifier.py      & 比对 QR 与 LLM 提取结果，返回校验结论 \\
重复检测器  & pipeline/duplicate\_checker.py & 按发票号码查重，记录入库 \\
\bottomrule
\end{longtable}

\subsection{处理流程图}

系统处理流程分为 5 个阶段，通过 SSE（Server-Sent Events）实时推送进度和中间结果：

\begin{enumerate}[label=阶段 \arabic*:, leftmargin=*]
  \item \textbf{YOLO 检测}：上传图片 → YOLO 推理 → 返回二维码裁剪区域和检测框坐标 → 推送检测结果（含 bbox 可视化预览）。
  \item \textbf{二维码解析}：pyzbar 解码裁剪区域 → 解析发票号码、金额、日期 → 推送解析结果。
  \item \textbf{OCR 文字识别}：PaddleOCR 对整图推理 → 返回文本行列表 → 推送 OCR 文字。
  \item \textbf{大模型提取}：图片 + OCR 文字 → 多模态 LLM → JSON 结构化输出 → 推送提取字段；若失败则降级规则提取。
  \item \textbf{双信息校验}：逐字段比对 QR 与 LLM 结果 → 输出一致/不一致 → 推送校验结果 → 等待用户确认后跳转。
\end{enumerate}

\begin{figure}[H]
\centering
\includegraphics[width=\textwidth]{screenshots/flowchart.drawio.png}
\caption{系统处理流程图}
\label{fig:flowchart}
\end{figure}


% ============================================================
%  第 4 章  详细实现
% ============================================================
\section{详细实现}

\subsection{环境配置}

系统基于 Python 3.10 + conda 虚拟环境构建，关键依赖如下：

\begin{table}[H]
\centering
\begin{tabular}{lll}
\toprule
\textbf{类别}     & \textbf{库}                        & \textbf{用途} \\
\midrule
深度学习  & PyTorch 2.5.1 + CUDA 12.4      & YOLO 推理 + PaddleOCR 加速 \\
目标检测  & Ultralytics (YOLOv8)            & 二维码区域定位 \\
OCR       & PaddleOCR + PaddlePaddle 3.2.0  & 离线文字识别 \\
大模型    & dashscope / openai / anthropic   & 多后端 LLM 字段提取 \\
二维码    & pyzbar                           & 二维码解码 \\
Web       & FastAPI + Jinja2 + uvicorn       & Web 服务与模板渲染 \\
数据库    & SQLAlchemy + SQLite              & 数据持久化 \\
导出      & openpyxl + reportlab             & Excel/PDF 生成 \\
图像      & OpenCV + PyMuPDF                 & 图像处理 + PDF 转 PNG \\
\bottomrule
\end{tabular}
\caption{关键依赖清单}
\label{tab:deps}
\end{table}

安装命令：\texttt{pip install -r requirements.txt}（共 122 个包）。

环境注意事项：
\begin{itemize}
  \item PaddlePaddle 必须锁定 3.2.0 版本（3.3.1 存在 oneDNN/PIR 兼容性 bug，Windows + RTX 4060 触发）。
  \item 需设置环境变量 \texttt{KMP\_DUPLICATE\_LIB\_OK=TRUE} 解决 PyTorch 与 PaddlePaddle 的 OpenMP 库冲突。
  \item 大模型 API Key 通过 \texttt{.env} 文件配置，\texttt{.gitignore} 已排除该文件。
\end{itemize}

\subsection{模型下载与加载}

\begin{itemize}
  \item \textbf{YOLO 模型}：基于 YOLOv8n.pt 预训练权重微调，最优权重保存为 \texttt{data/model\_weights/yolo\_qr.pt}。Ultralytics 首次运行时自动下载预训练权重。
  \item \textbf{PaddleOCR 模型}：PaddleOCR 首次调用时自动下载检测和识别推理模型到 \texttt{\textasciitilde/.paddleocr/} 目录。系统在 FastAPI 启动事件中执行 OCR 预热推理（一张 100×100 纯白图片），避免首次用户请求等待模型加载。
\end{itemize}

\subsection{YOLO 二维码检测模块}

YOLO 检测器封装在 YOLOQRDetector 类中，接口极简：输入图片 → 输出裁剪后的二维码区域和边界框坐标。核心逻辑：(1) 调用 YOLO 模型推理，置信度阈值 0.2（降低漏检）；(2) 取最高置信度的检测框；(3) 从原图裁剪二维码区域返回。

\subsection{二维码解析模块}

\texttt{decode\_qr()} 函数接收裁剪后的二维码图片，调用 pyzbar 解码二进制数据，再由格式解析函数自动判断纸质发票或电子发票格式，提取发票号码、金额、日期等字段。电子发票格式的判断规则：若逗号分隔后的第 4 个字段为 8 位以上纯数字，则为电子发票格式。

\subsection{PaddleOCR 离线识别模块}

\texttt{ocr\_image()} 函数封装了 PaddleOCR 推理和结果整理：输入图片最长边超过 1024px 时等比缩放以平衡速度与精度；推理结果按垂直坐标排序以保持原文阅读顺序；返回简化为（文本, 置信度）元组列表，隐藏 PaddleOCR 内部结构。

\subsection{大模型多后端提取器}

LLM 提取器采用适配器模式支持 4 种 API 后端：DashScope（通义千问 Qwen-VL）、OpenAI（GPT-4o）、Anthropic（Claude）、DeepSeek。通过环境变量 \texttt{LLM\_BACKEND} 切换。每个后端实现：(1) 图片转 Base64；(2) 构造含 OCR 文字和 JSON Schema 的提示词；(3) 发送 HTTP 请求；(4) 解析 JSON 响应并验证必填字段。调用失败时自动降级为规则提取。

\subsection{双信息校验模块}

\texttt{verify()} 函数接收 QR 解析结果和 LLM 提取结果，逐字段比对以下三项：

\begin{enumerate}[label=(\arabic*)]
  \item \textbf{发票号码}：精确字符串匹配。
  \item \textbf{开票日期}：经日期标准化后匹配——支持 YYYYMMDD、YYYY-MM-DD、YYYY年MM月DD日 三种格式统一为 YYYY-MM-DD。
  \item \textbf{金额}：数值匹配，容差 ±0.02 元。同时比对不含税金额和价税合计两个字段，任一命中即通过——避免 QR 存储总金额而 LLM 提取不含税金额导致的错配。
\end{enumerate}

三项全部通过则判定校验通过，否则返回不通过及每个不一致字段的 QR 值与提取值对比。

\subsection{重复检测实现}

\texttt{check\_duplicate()} 函数按发票号码查询数据库，已存在则返回原记录。提交时，查重使用流水线原始识别值（task\_data），而非表单提交值——防止用户修改表单字段绕过查重。

\subsection{SSE 实时推送数据流}

处理进度通过 SSE 协议实时推送，6 种事件类型：

\begin{table}[H]
\centering
\begin{tabular}{lp{4cm} p{6cm}}
\toprule
\textbf{事件} & \textbf{触发阶段} & \textbf{数据内容} \\
\midrule
progress      & 各阶段开始  & 当前阶段标识 + 进度消息 \\
yolo\_result  & YOLO 完成后 & 是否检测到、bbox 坐标、带框预览图 \\
qr\_result    & QR 解析后   & 是否解析成功、解析出的字段 \\
ocr\_result   & OCR 完成后  & 文本行列表（文本 + 置信度） \\
llm\_result   & LLM 完成后  & 提取来源（llm/rules）、JSON 字段 \\
verify\_result & 校验完成后 & 是否通过、二维码是否缺失、不一致列表 \\
\bottomrule
\end{tabular}
\caption{SSE 事件类型}
\label{tab:sse}
\end{table}

所有中间结果同时存入服务端 \texttt{app.state.tasks[task\_id]}，后续表单页和校验失败页从 task store 读取——数据流在服务端闭环，不依赖浏览器 sessionStorage。


% ============================================================
%  第 5 章  实验结果与展示
% ============================================================
\section{实验结果与展示}

\subsection{YOLO 模型训练结果}

训练数据集通过以下流程构建：百度图片爬取 → dHash 感知哈希去重 → OpenCV 初筛 → YOLO 引导式重标注 → 数据增强（6 种变换，bbox 同步更新）。

最终训练集规模：
\begin{itemize}
  \item 原始标注图片：123 张
  \item 数据增强后：1187 张（约 9.6 倍扩充）
  \item 训练集 / 验证集：949 张 / 238 张（80:20 分割）
  \item 检测类别：1 类（qr\_code）
\end{itemize}

训练配置：YOLOv8n，50 epochs，batch=8，imgsz=640。

训练结果：
\begin{itemize}
  \item mAP50：0.985
  \item mAP50-95：0.788
  \item 单张推理耗时：约 2.8ms（RTX 4060）
\end{itemize}

\begin{figure}[H]
\centering
\includegraphics[width=0.9\textwidth]{screenshots/yolo_training_results.png}
\caption{YOLO 训练曲线（results.png）}
\label{fig:yolo_train}
\end{figure}

\subsection{YOLO 检测效果}

在真实增值税发票上的检测效果：

\begin{figure}[H]
\centering
\includegraphics[width=0.7\textwidth]{screenshots/yolo_detection.png}
\caption{YOLO 二维码检测结果（绿色框为检测到的 QR 区域）}
\label{fig:yolo_detect}
\end{figure}

\subsection{系统功能演示}

以下截图展示完整的端到端使用流程：

\begin{figure}[H]
\centering
\includegraphics[width=0.85\textwidth]{screenshots/01_upload.png}
\caption{上传发票页面——支持拖拽或点击选择，JPG/PNG/PDF 格式}
\label{fig:upload}
\end{figure}

\begin{figure}[H]
\centering
\includegraphics[width=0.85\textwidth]{screenshots/02_process.png}
\caption{处理中页面——SSE 实时推送进度条、YOLO 检测、OCR 文字、LLM 提取结果}
\label{fig:process}
\end{figure}

\begin{figure}[H]
\centering
\includegraphics[width=0.85\textwidth]{screenshots/03_verify.png}
\caption{校验结果——双信息比对一致，显示"校验通过"，等待用户确认}
\label{fig:verify}
\end{figure}

\begin{figure}[H]
\centering
\includegraphics[width=0.85\textwidth]{screenshots/04_form.png}
\caption{报销表单——7 个发票字段自动填入，用户补填报销人/事由/部门}
\label{fig:form}
\end{figure}

\begin{figure}[H]
\centering
\includegraphics[width=0.85\textwidth]{screenshots/05_success.png}
\caption{提交成功——展示识别结果可视化图，提供 Excel/PDF 双格式下载}
\label{fig:success}
\end{figure}

\begin{figure}[H]
\centering
\includegraphics[width=0.85\textwidth]{screenshots/06_duplicate.png}
\caption{重复报销报警——同一发票号码再次提交时阻塞并展示原记录}
\label{fig:duplicate}
\end{figure}

\subsection{端到端测试}

编写了完整的端到端自动化测试脚本，覆盖上传 → SSE 全流程跟踪 → 结果汇总。测试覆盖 YOLO 检测、二维码解析、OCR 识别、LLM 提取、校验、填单、查重、导出的完整链路。34 个单元测试覆盖 6 个核心模块，全部通过。


% ============================================================
%  第 6 章  问题与解决方法
% ============================================================
\section{问题与解决方法}

开发过程中遇到并解决了 15 项技术问题，按类别汇总如下：

\subsection{环境与依赖兼容性}

\textbf{PaddlePaddle 版本兼容性}：PaddlePaddle 3.3.1 在 Windows + RTX 4060 环境触发 oneDNN/PIR bug，经多版本测试（3.0rc1、2.6.2）后锁定 PaddlePaddle 3.2.0 作为唯一兼容版本。

\textbf{PyTorch 与 PaddlePaddle 库冲突}：两个框架链接了不同版本的 OpenMP 库，需设置 \texttt{KMP\_DUPLICATE\_LIB\_OK=TRUE} 解决。

\textbf{PaddleOCR API 适配（3 轮）}：PaddleOCR 3.5 移除了 use\_gpu、show\_log 参数；PaddlePaddle 3.2.0 下 OCRResult 为字典式（rec\_texts / rec\_scores / rec\_polys），需按字典访问。

\subsection{Web 框架适配}

\textbf{Starlette 1.0 变更}：Jinja2 裸 Environment 不可直接用于 FastAPI，需使用 Jinja2Templates；TemplateResponse 的 request 参数从自动注入变为必须显式传入第一个位置参数——6 处调用全部修复。

\subsection{图像处理}

\textbf{cv2.imread 中文路径}：Windows 下 OpenCV 不支持中文路径，统一改用 np.frombuffer + cv2.imdecode 模式读取图片文件。

\textbf{PDF 处理}：process 页面硬编码了 jpg 扩展名，PDF 转 PNG 后路径不匹配；修复为上传时返回真实扩展名，前端从 URL 参数读取。

\textbf{OCR 性能}：PaddleOCR 处理大图耗时 14.6s，通过最长边 1024px 缩放降至 9.4s（下降 36\%）；通过 FastAPI 启动预热避免首次请求等待模型加载。

\subsection{业务逻辑}

\textbf{二维码格式错位}：电子发票 QR 编码格式与纸质发票不同，字段偏移 1 位——通过自动检测第 4 个字段是否为 8 位以上纯数字来判断格式。

\textbf{日期格式不匹配}：QR 存 YYYYMMDD，LLM 返回 YYYY-MM-DD——引入日期标准化函数统一格式后比对。

\textbf{金额字段错配}：QR 存价税合计，LLM 的 amount 是不含税金额——金额校验改为同时比对 amount 和 total 字段，任一命中即通过。

\textbf{重复检测绕过}：表单字段可编辑，用户改发票代码即可绕过组合查重——改为仅按发票号码查重，且查重使用流水线原始识别值而非表单提交值。

\subsection{下载与导出}

\textbf{中文文件名乱码}：Content-Disposition 含中文时 Starlette 用 latin-1 编码导致 500 错误——改用 filename*=UTF-8'' + URL 编码。

\textbf{IDM 拦截下载}：Internet Download Manager 拦截直链附件下载——改用 JavaScript fetch + Blob + 动态创建 a 元素触发下载，绕过下载管理器拦截。

\subsection{YOLO 模型}

\textbf{真发票漏检}：初始模型对真实发票 QR 检测率低，置信度阈值从 0.5 降至 0.2，手动标注 3 张真实发票加入训练集，增强后重训，mAP50 升至 0.985。

\textbf{极小 QR 未检出}：部分发票中 QR 码仅占图 8.8\% 面积，训练集中此类样本稀缺——已知限制，需更多极小 QR 训练样本。


% ============================================================
%  第 7 章  总结与改进
% ============================================================
\section{总结与改进}

\subsection{系统优点}

\begin{enumerate}[label=(\arabic*)]
  \item \textbf{端到端自动化}：从发票上传到报销单导出的全流程无人干预，大幅提升财务效率。
  \item \textbf{双信息校验}：二维码机器数据与 OCR+LLM 提取人工可读数据的交叉校验，是防止信息提取错误的有效机制。
  \item \textbf{降级设计}：YOLO 未检测到二维码时跳过校验继续流程；LLM 不可用时自动降级规则提取——不因单点故障阻塞整体流程。
  \item \textbf{模块化架构}：7 个 pipeline 模块可独立测试、独立替换，不依赖 Web 框架，符合工程设计原则。
  \item \textbf{实时反馈}：SSE 推送进度条 + 各阶段中间结果可视化，用户体验透明可控。
  \item \textbf{多后端支持}：LLM 适配器模式支持 4 种 API，不影响业务代码。
  \item \textbf{离线 OCR}：发票图片不出本地，满足数据隐私要求。
\end{enumerate}

\subsection{不足与未来优化方向}

\begin{enumerate}[label=(\arabic*)]
  \item \textbf{YOLO 模型泛化}：当前模型对极小 QR 码（占图不足 10\%）的检测率偏低，需收集更多不同分辨率和拍摄距离的训练样本。
  \item \textbf{发票类型扩展}：目前仅支持增值税发票，出租车票、定额发票、过路费票等不在范围内——扩展需要新增版面分析模块和字段定义。
  \item \textbf{批量处理}：当前一次仅处理一张发票，财务实际场景需要批量上传和批量处理能力。
  \item \textbf{多用户与权限}：系统缺乏用户认证和角色权限管理，不适用于多人员协作场景。
  \item \textbf{生产部署}：当前仅支持 uvicorn 单进程运行，需加入反向代理（Nginx）、进程管理（Gunicorn）、数据库升级（MySQL/PostgreSQL）等才能上线。
  \item \textbf{真伪验证}：本系统仅做信息提取和查重，未对接税务局接口进行发票真伪验证——这是财务合规的重要补充。
  \item \textbf{移动端适配}：当前为桌面 Web 设计，未做响应式移动端适配。
\end{enumerate}


% ============================================================
%  参考文献
% ============================================================
\clearpage
\phantomsection
\addcontentsline{toc}{section}{参考文献}
\begin{thebibliography}{99}

\bibitem{yolov8}
G. Jocher, A. Chaurasia, and J. Qiu, ``Ultralytics YOLOv8,'' 2023.
\url{https://github.com/ultralytics/ultralytics}

\bibitem{paddleocr}
Baidu Inc., ``PaddleOCR: Awesome Multilingual OCR Toolkits Based on PaddlePaddle,'' 2023.
\url{https://github.com/PaddlePaddle/PaddleOCR}

\bibitem{qwen-vl}
Alibaba Cloud, ``Qwen-VL: A Versatile Vision-Language Model,'' 2024.
\url{https://github.com/QwenLM/Qwen-VL}

\bibitem{pyzbar}
``pyzbar: Read One-Dimensional Barcodes and QR Codes from Python,''
\url{https://github.com/NaturalHistoryMuseum/pyzbar}

\bibitem{fastapi}
S. Ram\'irez, ``FastAPI: Modern Web Framework for Building APIs with Python,''
\url{https://fastapi.tiangolo.com/}

\bibitem{vat-invoice}
国家税务总局, ``增值税电子发票税控开票软件数据接口规范,'' 2019.

\bibitem{reportlab}
``ReportLab: Open-Source PDF Toolkit for Python,''
\url{https://www.reportlab.com/}

\bibitem{openpyxl}
``openpyxl: A Python Library to Read/Write Excel 2010 xlsx/xlsm Files,''
\url{https://openpyxl.readthedocs.io/}

\end{thebibliography}


% ============================================================
%  附录
% ============================================================
\appendix
\section{关键代码片段}

\subsection{YOLO 检测器核心代码}

\begin{lstlisting}[language=Python]
class YOLOQRDetector:
    def __init__(self):
        self.model = YOLO(YOLO_MODEL_PATH)

    def detect(self, image):
        results = self.model(image, conf=YOLO_CONFIDENCE, verbose=False)
        if not results or len(results[0].boxes) == 0:
            return None, None
        boxes = results[0].boxes
        best = boxes[boxes.conf.argmax()]
        x1, y1, x2, y2 = map(int, best.xyxy[0].tolist())
        cropped = image[y1:y2, x1:x2]
        return cropped, (x1, y1, x2, y2)
\end{lstlisting}

\subsection{校验器核心代码}

\begin{lstlisting}[language=Python]
def verify(qr: QRData, extracted: dict) -> VerificationResult:
    mismatches = []

    # invoice number matching
    if qr.invoice_number != extracted.get("invoice_number", ""):
        mismatches.append({"field": "invoice_number",
            "qr": qr.invoice_number,
            "extracted": extracted.get("invoice_number", "")})

    # date matching (after normalization)
    qr_date = _normalize_date(qr.date)
    ext_date = _normalize_date(extracted.get("invoice_date", ""))
    if qr_date != ext_date:
        mismatches.append({"field": "invoice_date",
            "qr": qr_date, "extracted": ext_date})

    # amount matching (check both amount and total, tolerance 0.02)
    qr_amt = float(qr.amount)
    ext_amt = extracted.get("amount")
    ext_tot = extracted.get("total")
    if not (_close(qr_amt, ext_amt) or _close(qr_amt, ext_tot)):
        mismatches.append({"field": "amount",
            "qr": qr_amt, "extracted": ext_amt or ext_tot})

    return VerificationResult(
        passed=len(mismatches) == 0,
        mismatches=mismatches)
\end{lstlisting}

\subsection{SSE 推送辅助函数}

\begin{lstlisting}[language=Python]
def _sse(event: str, data: dict) -> str:
    return f"event: {event}\n" \
           f"data: {json.dumps(data, ensure_ascii=False)}\n\n"
\end{lstlisting}

\end{document}
